% PAGES: 115-116
% NOTE: p.115 opens with the tail of the "One Period Binomial Model Example
% (Continued)" block begun at the end of sec03_c2_3.tex (PDF p.114). That
% block is typeset manually (bold-italic title + \rule), not as a LaTeX
% environment. The small square printed before the (3.41) tag closes the
% EXAMPLE, not a proof -- the preceding proof was already closed in
% sec03_c2_3.tex. Closed here with \hfill$\square$ + rule, matching the
% convention used by every other example block in this chapter
% (cf. sec03_c1_2.tex).

From (3.39) and (3.40), we find that the risk--neutral probabilities
associated to each of the two states of the market at time $\tau$ are
\begin{align*}
p_{RN}(1) &= e^{r\delta t} Q^1 = \frac{e^{r\delta t}-d}{u-d}; \\
p_{RN}(2) &= e^{r\delta t} Q^2 = \frac{u-e^{r\delta t}}{u-d}.
\end{align*}

Then, we obtain from (3.31) that the value at time $t_0$ of a derivative
security with payoffs at time $\tau$ equal to $V_\tau(1)$, if the ``up''
state occurs, and equal to $V_\tau(2)$, if the ``down'' state occurs, is
\begin{align}
V_{t_0} &= e^{-r\delta t}\left(p_{RN}(1)V_\tau(1) \ + \ p_{RN}(2)V_\tau(2)\right) \notag \\
        &= e^{-r\delta t}\left(\frac{e^{r\delta t}-d}{u-d}V_\tau(1) \ + \ \frac{u-e^{r\delta t}}{u-d}V_\tau(2)\right).
\end{align}
\hfill$\square$

\medskip
\noindent\rule{\linewidth}{0.4pt}
\medskip

We conclude this section by providing further intuition on the
no--arbitrage condition (3.11) for a one period market model, if the
market is also complete. From the definition (3.32) of the discrete
risk--neutral state probabilities at time $\tau$, we find that
\[
Q_k \ = \ \left(\sum_{i=1}^n Q_i\right) p_{RN}(k).
\]

Then, the no--arbitrage condition (3.11) can be written as
\begin{align}
S_{t_0} &= \sum_{k=1}^n Q_k S_\tau^{\omega^k} \ = \ \sum_{k=1}^n \left(\sum_{i=1}^n Q_i\right) p_{RN}(k) S_\tau^{\omega^k} \notag \\
        &= \left(\sum_{i=1}^n Q^i\right) \sum_{k=1}^n p_{RN}(k) S_\tau^{\omega^k} \notag \\
        &= e^{-r\delta t} \sum_{i=1}^n p_{RN}(k) S_\tau^{\omega^k},
\end{align}
where the last equality follows from (3.38).

Thus, in a complete market, the no--arbitrage condition (3.11) states that
the initial values at time $t_0$ of the $m$ securities are equal to the
present values of the risk--neutral expected values of the securities at
time $\tau$.

\subsection{State prices}

We can now explain why the values $Q^i > 0$, $i = 1:n$, which appeared in
the arbitrage--free condition (3.10) for a one period market model are
called state prices.

Consider a derivative security paying \$1 if state $\omega^p$ occurs at
time $\tau$ and 0 otherwise; this derivative security is also called an
insurance contract on the state $\omega^p$. In other words, if $C_\tau$
is the $n \times 1$ price vector of the insurance contract at time
$\tau$, then
\[
C_\tau^{\omega^k} = 0, \ \ \forall\, 1 \le k \neq p \le n; \qquad C_\tau^{\omega^p} = 1.
\]

If the market is complete and arbitrage--free, it follows from (3.30)
that the value $C_{t_0}$ of the insurance contract at time $t_0$ is
\[
C_{t_0} \ = \ \sum_{k=1}^n C_\tau^{\omega^k} Q_k \ = \ Q_p.
\]

We conclude that, for any $p = 1:n$, $Q_p$ is the price of the insurance
contract paying \$1 if state $p$ occurs, which can also be regarded as
the price at time $t_0$ of state $\omega^p$ occurring at time $\tau$.

\section{A one period index options market model}

A snapshot taken on March 9, 2012, of the mid prices (i.e., the average
of the bid and ask prices) and trading volumes of S\&P 500 options
maturing on December 22, 2012, (i.e., 12/22/2012), corresponding to a
spot price of the index of $1,370$, can be found in Table 3.1.

\begin{table}[H]
\centering
\caption{Dec 2012 SPX option prices on 3/9/2012}
\begin{tabular}{|c|c|c||c|c|c|}
\hline
Call Strike & Price & Volume & Put Strike & Price & Volume \\
\hline
C1175 & 225.40 & 250  & P1175 & 46.60  & 1     \\
\hline
C1200 & 205.55 & 215  & P1200 & 51.55  & 3204  \\
\hline
C1225 & 186.20 & 1    & P1225 & 57.15  & 1401  \\
\hline
C1250 & 167.50 & 650  & P1250 & 63.30  & 104   \\
\hline
C1275 & 149.15 & 163  & P1275 & 70.15  & 56    \\
\hline
C1300 & 131.70 & 1    & P1300 & 77.70  & 150   \\
\hline
C1325 & 115.25 & 40   & P1325 & 86.20  & 200   \\
\hline
C1350 & 99.55  & 320  & P1350 & 95.30  & 10118 \\
\hline
C1375 & 84.90  & 1002 & P1375 & 105.30 & 1250  \\
\hline
C1400 & 71.10  & 5300 & P1400 & 116.55 & 1250  \\
\hline
C1425 & 58.70  & 4    & P1425 & 129.00 & 200   \\
\hline
C1450 & 47.25  & 9050 & P1450 & 143.20 & 1     \\
\hline
C1500 & 29.25  & 1000 & P1500 & 173.95 & 6     \\
\hline
C1550 & 15.80  & 1000 & P1550 & 210.80 & 9     \\
\hline
C1575 & 11.10  & 200  & P1575 & 230.90 & 0     \\
\hline
C1600 & 7.90   & 546  & P1600 & 252.40 & 9     \\
\hline
\end{tabular}
\end{table}

We use a small number of these options to construct an arbitrage--free
complete one period market model for the 12/22/2012 S\&P options as
follows:\footnote{The author first learned of this type of one period
market models for options from the regretted Salih Neftci; see also
Chapter 11 from the monograph Neftci \cite{31}.} We choose one call
option and one put option with the same strike, with all the other
options having different strikes. The strikes of the options will cover
the majority of the range of strikes, and are best not to be close to
each other; e.g., the strikes of the options from (3.43) are at least 50
apart. We will subsequently show that the rest of the options from
Table 3.1 are accurately priced in this model.
